\section{Lời giải Bộ 31 (THPT Như Thanh, Thanh Hóa)}

\deda{
Trong mặt phẳng với hệ tọa độ $Oxy$ cho đường tròn $(C): x^2+y^2-2x+4y-20=0$ và đường thẳng $d:3x+4y-20=0$. Chứng minh rằng $d$ tiếp xúc với $(C)$. Tam giác $ABC$ có đỉnh $A\in (C)$, hai đỉnh $B,C\in d$, trung điểm cạnh $AB\in (C)$. Tìm tọa độ các đỉnh tam giác $ABC$, biết trực tâm tam giác $ABC$ trùng với tâm đường tròn $(C)$ và $B$ có hoành độ dương.
}
\lgi

\definecolor{uuuuuu}{rgb}{0.26666666666666666,0.26666666666666666,0.26666666666666666}
\definecolor{qqqqff}{rgb}{0.,0.,1.}
\begin{tikzpicture}[line cap=round,line join=round,>=triangle 45,x=1.0cm,y=1.0cm]
\clip(-4.3,-2.44) rectangle (11.32,6.3);
\draw (1.76,4.12)-- (0.56,0.02);
\draw (4.64,2.18)-- (1.76,4.12);
\draw (5.032111042372097,2.9574974092966118) node[anchor=north west] {(d)};
\draw(2.447300509337861,2.821765704584041) circle (1.4907361523523388cm);
\draw (1.76,4.12)-- (3.192702702702703,1.4137837837837839);
\draw (1.4642410958904108,3.109490410958904)-- (4.64,2.18);
\draw (1.16,2.07)-- (3.7346010186757224,3.5735314091680825);
\draw (2.8430884082715373,4.686641338754195) node[anchor=north west] {(C)};
\draw (0.56,0.02)-- (5.68,2.72);
\begin{scriptsize}
\draw [fill=qqqqff] (1.76,4.12) circle (1.5pt);
\draw[color=qqqqff] (1.88654070261415,4.373923819596973) node {$A$};
\draw [fill=qqqqff] (0.56,0.02) circle (1.5pt);
\draw[color=qqqqff] (0.359743403199474,0.2534106259959226) node {$B$};
\draw [fill=qqqqff] (4.64,2.18) circle (1.5pt);
\draw[color=qqqqff] (4.774578967772031,2.4424332600964807) node {$C$};
\draw [fill=uuuuuu] (1.16,2.07) circle (1.5pt);
\draw[color=uuuuuu] (0.9667832933282007,2.3504575191678856) node {$F$};
\draw [fill=uuuuuu] (3.192702702702703,1.4137837837837839) circle (1.5pt);
\draw[color=uuuuuu] (3.119015631057322,1.0995874425389955) node {$A'$};
\draw [fill=uuuuuu] (2.447300509337861,2.821765704584041) circle (1.5pt);
\draw[color=uuuuuu] (2.548766037300034,3.2150294838966778) node {$H$};
\draw [fill=uuuuuu] (1.4642410958904108,3.109490410958904) circle (1.5pt);
\draw[color=uuuuuu] (1.5002425907140513,2.9574974092966118) node {$K$};
\draw [fill=uuuuuu] (3.7346010186757224,3.5735314091680825) circle (1.5pt);
\draw[color=uuuuuu] (3.7444506693717683,3.89564996676828) node {$E$};
\end{scriptsize}
\end{tikzpicture}


\deda{
Giải phương trình: $4\sqrt{5x^3-6x^2+2}+4\sqrt{-10x^3+8x^2+7x-1}+x-13=0$.
}
\lgi
ĐK: $\begin{cases} 5x^3-6x^2+2\ge 0\\ -10x^3+8x^2+7x-1\ge 0  \end{cases}$. $(*)$\\
\begin{eqnarray*}
PT&\iff 4\sqrt{5x^3-6x^2+2}+4\sqrt{-10x^3+8x^2+7x-1}=x-13\\
\end{eqnarray*}
Với $(*)$, áp dụng bất đẳng thức Cô-si ta được:\\
\cl{$4\sqrt{5x^3-6x^2+2}=4\sqrt{1.(5x^3-6x^2+2)}\le 4.\dfrac{5x^3-6x^2+2+1}{2}$.$(1)$}
Tương tự ta có\\
\cl{$4\sqrt{-10x^3+8x^2+7x-1}=2\sqrt{4(-10x^3+8x^2+7x-1)}\le 2.\dfrac{-10x^3+8x^2+7x-1+4}{2}$.$(2)$}
Từ $(1)$ và $(2)$ ta có:\\
\cl{$VT=4\sqrt{5x^3-6x^2+2}+4\sqrt{-10x^3+8x^2+7x-1}\le -4x^2+7x+9$.}
Mặt khác ta lại có:\\
\cl{$-4x^2+7x+9\le -4x^2+7x+9+4(x-1)^2=x-13=VP$.}
Vậy phương trình đã cho $\iff \begin{cases} 5x^3-6x^2+2=1\\    -10x^3+8x^2+7x-1=4 \\ 4(x-1)^2=0 \end{cases}\iff x=1$.\\
Vậy phương trình đã cho nghiệm duy nhất $x=1$.


\deda{
Cho các số $a,b,c\in R, a^2+b^2+c^2\neq 0$ và $2(4a^2+4b^2+c^2)=(2a+2b+c)^2$.\\
Tìm giá trị lớn nhất, nhỏ nhất của biểu thức $A=\dfrac{8a^3+8b^3+c^3}{(2a+2b+2c)(4ab+2bc+2ca)}$.
}
\lgi
Ta có:\\
\cl{$GT\iff \left[a+b+\dfrac{c}{2}\right]^2=2\left[a^2+b^2+\left(\dfrac{c}{2}\right)^2\right]$.}
Đặt $d=\dfrac{c}{2}$, ta có:\\
\cl{$(a+b+d)^2=2(a^2+b^2+d^2)$. $(1)$}
Mặt khác ta có:\\
\cl{$ab+bd+da=\dfrac{1}{2}\left[(a+b+d)^2-a^2-b^2-d^2\right]$. $(2)$}
Từ $(1)$ và $(2)$ ta được $ab+bd+da=\dfrac{1}{2}(a^2+b^2+d^2)$.
Khi đó ta có:\\
$A=\dfrac{8a^3+8b^3+c^3}{2(2a+2b+c)(2ab+bc+ca)}
=\dfrac{a^3+b^3+\left(\dfrac{c}{2}\right)^3}{\left(a+b+\dfrac{c}{2}\right)\left(ab+\dfrac{bc}{2}+\dfrac{ca}{2}\right)}$\\
$=\dfrac{a^3+b^3+d^3}{(a+b+d)(ab+bd+da)}=\dfrac{4(a^3+b^3+d^3)}{(a+b+d)^3}=\dfrac{1}{16}\left[ \left( \dfrac{4a}{a+b+d} \right)^3+\left(\dfrac{4b}{a+b+d}\right)^3+\left(\dfrac{4d}{a+b+d}\right)^3 \right]$.\\
Đặt $x=\dfrac{4a}{a+b+d},y=\dfrac{4b}{a+b+d},z=\dfrac{4d}{a+b+d}$, khi đó ta có:\\
$\begin{cases}
x+y+z=4 \\ xy+yz+zx=4
\end{cases}\iff \begin{cases} y+z=4-x \\ yz=x^2-4x+4 \end{cases}$, với điều kiện $(y+z)^2\ge 4yz\Rightarrow 0\le x\le \dfrac{8}{3}$.

$\Rightarrow A=\dfrac{1}{16}\left [ x^3+(y+z)^3-3yz(y+z) \right ]=\dfrac{1}{16}\left [ 3x^3-12x^2+12x+16 \right ]$.\\

Xét hàm $f(x)=\dfrac{1}{16}\left [ 3x^3-12x^2+12x+16 \right ]$ với $0\le x\le \dfrac{8}{3}$.\\

$MinA=\dfrac{1}{16}$ khi $Minf(x)=1$, đặt được khi $a=0,c=2b\neq 0$.\\
$MaxA=\dfrac{1}{16}$ khi $Maxf(x)=\dfrac{11}{9}$, đạt được khi $a=b,c=8a,a\neq 0$.

\section{Lời giải Bộ 32 (THPT Chuyên Hạ Long)}

\deda{
Trong mặt phẳng với hệ trục tọa độ $Oxy$, cho tam giác $ABC$ có $A(2;6),B(1;1),C(6;3)$.\\
1. Viết phương trình đường tròn ngoại tiếp tam giác $ABC$.\\
2. Tìm trên các cạnh $AB,BC,CA$ các điểm $K,H,I$ sao cho chu vi tam giác $KHI$ nhỏ nhất.
}
\lgi
\vspace{5cm}

\deda{
Giải hệ phương trình \\ \cl{$\begin{cases} 3y\sqrt{2+x}+8\sqrt{2+x}=10y-3xy+12\\ 5y^3\sqrt{2-x}-8=6y^2+xy^3\sqrt{2-x} \end{cases}$.} 
}
\lgi
ĐK: $x\in [-2;2]$\\
Nhận xét $y=0$ không thỏa mãn phương trình $(2)$\\
$(2)\iff \left ( \sqrt{2-x}\right )^3+3\sqrt{2-x}=\left(\dfrac{2}{y}\right)^3+3.\dfrac{2}{y}$ $(*)$\\
Xét hàm số $f(t)=t^3+3t$ trên $R$ ta thấy hàm số đồng biến trên $R$.\\
$(*)\iff f\left(\sqrt{2-x}\right)=f\left(\dfrac{2}{y}\right)\iff \sqrt{2-x}=\dfrac{2}{y}$.\\
Thế vào PT $(1)$ ta có\\
$(1)\iff 3y\sqrt{2+x}+8\sqrt{2+x}=10y-3xy+12$\\
$\iff 3\sqrt{2+x}+4\sqrt{2+x}\sqrt{2-x}=10-3x+6\sqrt{2-x}$\\
$\iff 3\sqrt{2+x}-6\sqrt{2-x}+4\sqrt{4-x^2}+3x-10=0~(**)$\\

Đặt $\sqrt{2+x}-2\sqrt{2-x}=t\Rightarrow t^2=10-3x-4\sqrt{4-x^2}$\\
Phương trình $(**)$ trở thành $3t-t^2=0\iff t=0$ hoặc $t=3$\\
+ Với $t=0$ ta được $x=\dfrac{6}{5},y=\sqrt{5}$\\
+ Với $t=3$ ta có $\sqrt{2+x}-2\sqrt{2-x}=3$, phương trình vô nghiệm vì $VT\le 2$.

\deda{
Chứng minh rằng: Với mọi $\Delta ABC$ ta đều có\\
\cl{$\left ( \sin \dfrac{A}{2}+\sin \dfrac{B}{2}+ \sin \dfrac{C}{2} \right )\left ( \cot \dfrac{A}{2}+\cot \dfrac{B}{2}+ \cot \dfrac{C}{2} \right )\ge \dfrac{9\sqrt{3}}{2}$.} 
}
\lgi
Ta có: $\dfrac{A}{2},\dfrac{B}{2},\dfrac{C}{2}\in \left(0,\dfrac{\pi}{2}\right)\Rightarrow \sin \dfrac{A}{2},\sin \dfrac{B}{2}.\sin \dfrac{C}{2},\cos \dfrac{A}{2},\cos \dfrac{B}{2},\cos \dfrac{C}{2}>0$\\
Có: $\sin \dfrac{A}{2}+\sin \dfrac{B}{2}+\sin \dfrac{C}{2}\ge 3\sqrt[3]{\sin \dfrac{A}{2}\sin \dfrac{B}{2}\sin \dfrac{C}{2}}> 0$ (Cauchy)\\
Lại có:
$\cot \dfrac{A}{2}+\cot \dfrac{B}{2}+\cot \dfrac{C}{2}$\\
$=\dfrac{\sin \dfrac{A}{2}\left(\sin \dfrac{B}{2}\cos \dfrac{C}{2}+\sin \dfrac{C}{2}\cos \dfrac{B}{2}\right)}{2\sin \dfrac{A}{2}\sin \dfrac{B}{2}\sin \dfrac{C}{2}}+\dfrac{\sin \dfrac{B}{2}\left(\sin \dfrac{A}{2}\cos \dfrac{C}{2}+\sin \dfrac{C}{2}\cos \dfrac{A}{2}\right)}{2\sin \dfrac{A}{2}\sin \dfrac{B}{2}\sin \dfrac{C}{2}}+\dfrac{\sin \dfrac{C}{2}\left(\sin \dfrac{A}{2}\cos \dfrac{B}{2}+\sin \dfrac{B}{2}\cos \dfrac{A}{2}\right)}{2\sin \dfrac{A}{2}\sin \dfrac{B}{2}\sin \dfrac{C}{2}}$\\
$=\dfrac{\sin \dfrac{A}{2}\cos \dfrac{A}{2}+\sin \dfrac{B}{2}\cos \dfrac{B}{2}+\sin \dfrac{C}{2}\cos \dfrac{C}{2}}{2\sin \dfrac{A}{2}\sin \dfrac{B}{2}\sin \dfrac{C}{2}}$\\
$\ge \dfrac{3\sqrt[3]{\sin \dfrac{A}{2}\cos \dfrac{A}{2}\sin \dfrac{B}{2}\cos \dfrac{B}{2}\sin \dfrac{C}{2}\cos \dfrac{C}{2}}}{2\sin \dfrac{A}{2}\sin \dfrac{B}{2}\sin \dfrac{C}{2}}$ (Cauchy)\\
Vậy ta có $\left ( \sin \dfrac{A}{2}+\sin \dfrac{B}{2}+ \sin \dfrac{C}{2} \right )\left ( \cot \dfrac{A}{2}+\cot \dfrac{B}{2}+ \cot \dfrac{C}{2} \right )\ge \dfrac{9}{2}\sqrt[3]{\cot \dfrac{A}{2}\cot \dfrac{B}{2}\cot \dfrac{C}{2}}$\\
Mà $\cot \dfrac{A}{2}\cot \dfrac{B}{2}\cot \dfrac{C}{2}\ge 3\sqrt{3}$\\
Vậy ta có đpcm.


\section{Lời giải Bộ 33 (THPT chuyên Vĩnh Phúc, Vĩnh Phúc - Khối AB)}

\deda{
Trong mặt phẳng với hệ trục tọa độ $Oxy$, cho tam giác cân $ABC$ có đáy $BC$ nằm trên đường thẳng $d: 2x-5y+1=0$, cạnh $AB$ nằm trên đường thẳng $d':12x-y-23=0$. Viết phương trình đường thẳng $AC$ biết nó đi qua điểm $M(3;1)$.
}
\lgi 
\vspace{5cm}

\deda{
Giải hệ phương trình \\ \cl{$\begin{cases} \sqrt{5x^2+2xy+2y^2}+\sqrt{2x^2+2xy+5y^2}=3(x+y)\\ \sqrt{x+2y+1}+2\sqrt[3]{12x+7y+8}=2xy+x+5 \end{cases}$.}
}
\lgi
ĐK: $\begin{cases} 5x^2+2xy+2y^2\ge 0\\ 2x^2+2xy+5y^2\ge 0\\
x+2y+1\ge 0 \end{cases}\iff x+2y+1\ge 0$\\
Khi hệ có nghiệm, từ PT $(1)$ suy ra $x+y\ge 0$\\
Ta thấy: $\sqrt{5x^2+2xy+2y^2}=\sqrt{(x-y)^2+(2x+y)^2}\ge 2x+y$ $(*)$, dấu bằng xảy ra khi $x=y$\\
Tương tự $\sqrt{2x^2+2xy+5y^2}\ge x+2y$ $(**)$, dấu bằng xảy ra khi $x=y$.\\
Từ $(*)$ và $(**)$ ta có
$VT(1)=\sqrt{5x^2+2xy+2y^2}+\sqrt{2x^2+2xy+5y^2}\ge 3(x+y)=VP(1)$\\
Dấu đẳng thức xảy ra khi $x=y$\\
Thế $x=y$ vào PT $(2)$ ta được:\\
$\sqrt{3x+1}+2\sqrt[3]{19x+8}=2x^2+x+5$ điều kiện $x\ge -\dfrac{1}{3}$\\
$\iff 2(x^2-x)+\left ( x+1-\sqrt{3x+1}\right )+2\left(x+2-\sqrt[3]{19x+8}\right)=0$\\
$\iff (x^2-x)\left[2+\dfrac{1}{x+1+\sqrt{3x+1}}+2.\dfrac{1}{(x+2)^2+(x+2)\sqrt[3]{19x+8}+\sqrt[3]{(19x+8)^2}}\right]=0$\\
$\iff x^2-x=0$ (Biểu thức trong ngoặc to lớn hơn $0$)\\
$\iff \hay{x=0\Rightarrow y=0 \\ x=1\Rightarrow y=1}$ (thỏa mãn ĐK)\\
Vậy hệ phương trình có hai nghiệm $(x;y)=(0;0),(1;1)$. 

\deda{
Cho ba số thực dương $a,b,c$ thỏa mãn $a^2+b^2+c^2=3$.\\
Tìm giá trị nhỏ nhất của biểu thức $S=8(a+b+c)+5\left( \dfrac{1}{a}+\dfrac{1}{b}+\dfrac{1}{c}\right)$.
}
\lgi 
Ta sẽ chứng minh $8a+\dfrac{5}{a}\ge \dfrac{3a^2+23}{2}(1)$ với $0<a<\sqrt{3}$, dấu bằng khi $a=1$. Thật vậy\\
$8a+\dfrac{5}{a}\ge \dfrac{3a^2+23}{2}\iff (a-1)^2(3a-10)\le 0$ (Luôn đúng) \\ 
với mọi $0<a<\sqrt{3}$ dấu bằng khi $a=1$.\\
Tương tự\\
$8b+\dfrac{5}{b}\ge \dfrac{3b^2+23}{2}(2)$, dấu bằng khi $b=1$.\\
$8c+\dfrac{5}{c}\ge \dfrac{3c^2+23}{2}(3)$, dấu bằng khi $c=1$.\\
Từ $(1),(2),(3)$ ta có:\\ $S=8(a+b+c)+5\left(\dfrac{1}{a}+\dfrac{1}{b}+\dfrac{1}{c}\right)\ge \dfrac{3(a^2+b^2+c^2)+69}{2}=39$\\
Dấu bằng xảy ra khi $a=b=c=1$.
Vậy giá trị nhỏ nhất của $S=39$ đạt được  khi và chỉ khi $a=b=c=1$.\\
\textbf{Chú ý:} Để tìm ra vế phải của $(1)$ ta sử dụng phương pháp tiếp tuyến.\\

\section{Lời giải Bộ 34 (THPT chuyên Vĩnh Phúc, Vĩnh Phúc - Khối D)}

\deda{
Trong mặt phẳng với hệ trục tọa độ $Oxy$, cho đường tròn $(T): x^2+y^2-x-9y+18=0$ và hai điểm $A(4;1),B(3;-1)$. Gọi $C,D$ là hai điểm thuộc $(T)$ sao cho $ABCD$ là một hình bình hành. Viết phương trình đường thẳng $CD$.
}
\lgi
\vspace{5cm}

\deda{
Giải hệ phương trình \\ \cl{$\begin{cases} x^3-y^3+6x-3y=3x^2+4\\ x^2+6y+19=2\sqrt{3x+4}+3\sqrt{5y+14} \end{cases}$.} 
}
\lgi 
PT $(1)\iff (x-1)^3+3(x-1)=y^3+3y$\\
$\iff (x-1-y)\left[(x-1)^2+(x-1)y+y^2+3\right]=0$\\
$\iff x-1-y=0$ (Biểu thức trong ngoặc lớn hơn $0$)\\
$\iff y=x-1$\\
Thế $y=x-1$ vào PT $(2)$ ta được:\\
$x^2+6x+13=2\sqrt{3x+4}+3\sqrt{5x+9}~$ĐK: $x\ge -\dfrac{4}{3}$\\
$\iff (x^2+x)+2\left[(x+2)-\sqrt{3x+4}\right]+3\left[(x+3)-\sqrt{5x+9}\right]=0$\\
$\iff (x^2+x)\left[1+\dfrac{2}{x+2+\sqrt{3x+4}}+\dfrac{3}{x+3+\sqrt{5x+9}}\right]=0$\\
$\iff x^2+x=0$ (Biểu thức trong ngoặc lớn hơn $0$)\\
$\iff \hay{x=0\Rightarrow y=-1 \\ x=-1\Rightarrow y=-2}$(thỏa mãn ĐK)\\
Vậy hệ phương trình có hai nghiệm $(x;y)=(0;-1),(-1;-2)$
\deda{
Cho $a,b,c$ là các số thực dương thỏa mãn $a^2+b^2+c^2=1$. Chứng minh bất đẳng thức:\\ \centerline{$\left(\dfrac{1}{a}+\dfrac{1}{b}+\dfrac{1}{c}\right)-(a+b+c)\ge 2\sqrt{3}$.} 
}
\lgi 
Ta sẽ chứng minh \\
$-a+\dfrac{1}{a}\ge -2\sqrt{3}a^2+\dfrac{4\sqrt{3}}{3}$ $(1)$ với mọi $0<a<1$, dấu bằng khi $a=\dfrac{\sqrt{3}}{3}$\\
Thật vậy $(1)\iff \left(\sqrt{3}a-1\right)^2\left(2a+\sqrt{3}\right)\ge 0$ (Luôn đúng)\\
với mọi $0<a<1$, dấu bằng khi $a=\dfrac{\sqrt{3}}{3}$.\\
Tương tự\\
$-b+\dfrac{1}{b}\ge -2\sqrt{3}b^2+\dfrac{4\sqrt{3}}{3}$ $(2)$, dấu bằng khi $b=\dfrac{\sqrt{3}}{3}$\\
$-c+\dfrac{1}{c}\ge -2\sqrt{3}c^2+\dfrac{4\sqrt{3}}{3}$ $(3)$, dấu bằng khi $c=\dfrac{\sqrt{3}}{3}$\\
Từ $(1),(2),(3)$ ta có\\
$-(a+b+c)+\left (\dfrac{1}{a}+\dfrac{1}{b}+\dfrac{1}{c}\right)\ge -2\sqrt{3}(a^2+b^2+c^2)+3.\dfrac{4\sqrt{3}}{3}$\\
$\iff \left (\dfrac{1}{a}+\dfrac{1}{b}+\dfrac{1}{c}\right)-(a+b+c)\ge 2\sqrt{3}$\\
Dấu bằng xảy ra khi $a=b=c=\dfrac{\sqrt{3}}{3}$.\\
\textbf{Chú ý:} Để tìm ra vế phải của $(1)$ ta sử dụng phương pháp tiếp tuyến.

\section{Lời giải Bộ 35 (THPT Hồng Quang, Hải Dương)}

\deda{
Trong mặt phẳng tọa độ $Oxy$, cho tam giác $ABC$ có đường cao và đường trung tuyến kẻ từ đỉnh $A$ lần lượt có phương trình là $x-3y=0$ và $x+5y=0$. Đỉnh $C$ nằm trên đường thawgnr $\Delta: x+y-2=0$ và có hoành độ dương. Tìm tọa độ các đỉnh của tam giác $ABC$, biết đường thẳng chứa trung tuyến kẻ từ $C$ đi qua điểm $E(-2,6)$.
}
\lgi 
\vspace{5cm}

\deda{Giải hệ phương trình \\ \cl{$\begin{cases} x-\dfrac{1}{(x+1)^2}=\dfrac{y}{x+1}-\dfrac{1+y}{y}\\ \sqrt{8y+9}=(x+1)\sqrt{y}+2 \end{cases}$.} 
}
\lgi 
ĐK: $x\neq -1,y>0$\\
$x-\dfrac{1}{(x+1)^2}=\dfrac{y}{x+1}-\dfrac{1+y}{y}$\\
$\iff x+\dfrac{1+y}{y}=\dfrac{y}{x+1}+\dfrac{1}{(x+1)^2}$\\
$\iff \dfrac{xy+y+1}{y}-\dfrac{xy+y+1}{(x+1)^2}=0$\\
$\iff \hay{xy+y+1=0 \\ y=(x+1)^2}$\\
$*$ Với $y=(x+1)^2$, thay vào PT $(2)$ ta có $\sqrt{8(x+1)^2+9}=(x+1)|x+1|+2$\\
+ Xét $x>-1$, đặt $t=x+1$, $(t>0)$. Ta có phương trình\\
$\sqrt{8t^2+9}=t^2+2$\\
$\iff t^4-4t^2-5=0$\\
$\iff \hay{t^2=-1 \\ t^2=5}$\\
$\iff t=\sqrt{5}$ (Do $t>0$)\\
$\Rightarrow x=-1+\sqrt{5}\Rightarrow y=5$\\
+ Xét $x<-1$, đặt $t=x+1$, $(t<0)$. Ta có phương trình\\
$\sqrt{8t^2+9}=-t^2+2$\\
$\iff \va{t^4-12t^2-5=0 \\ t^2\le 2}$\\
Hệ vô nghiệm\\
$*$ Với $xy+y+1=0$, thay vào PT $(2)$ có $\sqrt{8y+9}+\dfrac{1}{y}\sqrt{y}-2=0$\\
Vì $y>0\Rightarrow \sqrt{8y+9}>3$\\
Phương trình vô nghiệm\\
Vậy hệ phương trình có nghiệm $\va{x=-1+\sqrt{5} \\ y=5}$.

\deda{
Cho các số dương $x,y,z$ thỏa mãn $x>y$ và $(x+z)(y+z)=1$.\\
Tìm giá trị nhỏ nhất của biểu thức
$P=\dfrac{1}{(x-y)^2}+\dfrac{4}{(x+z)^2}+\dfrac{4}{(y+z)^2}.$
}
\lgi 
Đặt $x+z=a$. Từ giả thiết ta có $y+z=\dfrac{1}{a}$.\\
Do $x>y\Rightarrow x+z>y+z\Rightarrow a>1$\\
Ta có $x-y=x+z-(y+z)=a-\dfrac{1}{a}=\dfrac{a^2-1}{a}$\\
$P=\dfrac{a^2}{(a^2-1)^2}+\dfrac{4}{a^2}+4a^2=\dfrac{a^2}{(a^2-1)^2}+3a^2+a^2+\dfrac{4}{a^2}\ge \dfrac{a^2}{(a^2-1)^2}+3a^2+4$ (Cauchy)\\
Đặt $t=a^2>1$. Xét hàm số $f(t)=\dfrac{t}{(t-1)^2}+3t+4$ với $t>1$\\
Ta có $f'(t)=\dfrac{-t-1}{(t-1)^3}+3\Rightarrow f'(t)=0\iff t=2$\\
Bảng biến thiên ...\\
a a a a a a\\
a a a a a a\\
a a a a a a\\
Từ bảng biến thiên có  $f(t)\ge 12\forall t>1$. Từ $(1)$ và $(2)$ ta có $P\ge 12$. Dấu bằng xảy ra\\
$\iff \va{x+z=\sqrt{2} \\ y+z=\dfrac{1}{\sqrt{2}}}$ Khi chẳng hạn $\va{x=1;z=\sqrt{2}-1 \\ y=\dfrac{1}{\sqrt{2}}-\sqrt{2}+1=1-\dfrac{1}{\sqrt{2}}}$\\
Vậy giá trị nhỏ nhất của $P$ là $12$.

\section{Lời giải Bộ 36 (THPT Lương Thế Vinh, Hà Nội) }

\deda{
Trong mặt phẳng tọa độ $Oxy$, cho hình chữ nhật $ABCD$ có diện tích bằng $15$. Đường thẳng $AB$ có phương trình $x-2y=0$. Trọng tâm của tam giác $BCD$ là điểm $G\left (\dfrac{16}{3};\dfrac{13}{3}\right)$. Tìm tọa độ bốn đỉnh của hình chữ nhật biết điểm $B$ có tung độ lớn hơn $3$.
}
\lgi 
\vspace{5cm}

\deda{
Giải hệ phương trình \\ \cl{$\begin{cases} 2x^3-3+2\sqrt{y^2+3y}=2x\sqrt{y}+y\\ x^2-\sqrt{y+3}+\sqrt{y}=0 \end{cases}$.} 
}
\lgi 
ĐK: $y\ge 0$\\
PT $(1)\iff 2x^3-2x\sqrt{y}+y=\left(\sqrt{y+3}-\sqrt{y}\right)^2=x^4$ (Do PT $(2)$)\\
$\iff \left(x^2-2x+\sqrt{y}\right)\left(x^2-\sqrt{y}\right)=0$\\
$*$ Với $\sqrt{y}=x^2$, PT $(2)\iff \sqrt{x^2+3}=2x^2$\\
$\iff 4x^4-x^2-3=0\iff x^2=1\Rightarrow (x;y)=(1;1),(-1;1)$\\
$*$ Với $\sqrt{y}=2x-x^2$ $(*)$, PT $(2)\iff \sqrt{3+(2x-x^2)^2}=2x$\\
$\iff \va{x\ge 0 \\ x^4-4x^3+3=0}\iff \hay{x=1\Rightarrow y=1 \\ x^2(x-3)-3x-3=0~(**)}$\\
Từ $(*)\Rightarrow 0\le x\le 2\Rightarrow (**)$ vô nghiệm.\\
Vậy hệ phương trình có nghiệm $(x;y)=(1;1),(-1;1)$
\deda{
Cho các số thực $a,b$ không âm và thỏa mãn: $3(a+b)+2(ab+1)\ge 5(a^2+b^2)$.\\
Tìm giá trị lớn nhất của biểu thức $T=3\sqrt{a+b}-3(a^2+b^2)+2(a+b)-ab$.
}
\hda 
Từ giả thiết ta có: $3(a+b)+2\ge 2(a+b)^2+3(a-b)^2\ge 2(a+b)^2~\forall a,b$\\
Đặt $t=a+b\ge 0\Rightarrow 0\le t\le 2$.
Ta có:\\
$T=ab+3\sqrt{a+b}-2(a^2+b^2)+1-(a+b-1)^2\le \left(\dfrac{a+b}{2}\right)^2+3\sqrt{a+b}-(a+b)^2+1$\\
$\Rightarrow T\le -\dfrac{3}{4}t^2+e\sqrt{t}+1=f(t), t\in [0;2]$\\
Ta có: $f'(t)=-\dfrac{3}{2}.\dfrac{t\sqrt{t}-1}{\sqrt{t}}$.
$f'(t)=0\iff t=1$\\
$f(0)=1;f(1)=\dfrac{13}{4};f(2)=3\sqrt{2}-2$.
Từ đó $Max T=\dfrac{13}{4}\iff t=1\iff a=b=\dfrac{1}{2}$.

\section{Lời giải Bộ 37 (THPT Nông Cống, Hà Nội) }

\deda{
Trong mặt phẳng với hệ tọa độ $Oxy$, cho đường tròn $(C)$ tâm $I$ $(x_I>0)$, $(C)$ đi qua điểm $A(-2;3)$ và tiếp xúc với đường thẳng $(d_1):x+y+4=0$ tại điểm $B$. $(C)$ cắt $(d_2):3x+4y-16=0$ tại $C$ và $D$ sao cho $ABCD$ là hình thang có hai đáy là $AD$ và $BC$, hai đường chéo $AC,BD$ vuông góc với nhau. Tìm tọa độ các điểm $B,C,D$.
}
\lgi 
\vspace{5cm}

\deda{
Giải hệ phương trình \\ \cl{$\begin{cases} \sqrt{x^2+xy+2y^2}+\sqrt{y^2+xy+2x^2}=2(x+y)\\ (8y-6)\sqrt{x-1}=\left(2+\sqrt{y-2}\right)\left(y+4\sqrt{x-2}+3\right)  \end{cases}$.} 
}
\hda 
ĐK: $x;y\ge 2$\\
PT $(1)\iff \sqrt{t^2+t+2}+\sqrt{2t^2+t+1}=2(t+1)$ với $t=\dfrac{x}{y},t>0$\\
Bình phương hai vế ta được $x=y$\\
Thay $x=y$ vào PT $(2)$ ta được:\\
$(8x-6)\sqrt{x-1}=\left(2+\sqrt{x-2}\right)\left(x+4\sqrt{x-2}+3\right)$\\
$\iff \sqrt{4x-4}\left[\left(\sqrt{4x-4}\right)^2+1\right]=\left(2+\sqrt{x-2}\right)\left[\left(2+\sqrt{x-2}\right)^2+1\right]$\\
Xét hàm số $f(t)=t^3+t$ luôn đồng biến trên $R$ nên $\sqrt{4x-4}=2+\sqrt{x-2}$\\
Giải ra ta được $\hay{x=2 \\ x=\dfrac{34}{9}}$\\
Vậy hệ có 2 nghiệm $(x;y)=(2;2),\left(\dfrac{34}{9};\dfrac{34}{9}\right)$.

\deda{
Cho $x,y$ là các số thực không âm thỏa mãn:\\
\cl{$\sqrt{2x^2+3xy+4y^2}+\sqrt{2y^2+3xy+4x^2}-3(x+y)^2\le 0$.}
Tìm giá trị nhỏ nhất của:\\
\cl{$P=2(x^3+y^3)+2(x^2+y^2)-xy+\sqrt{x^2+1}+\sqrt{y^2+1}$.}
}
\hda 
Ta có: $\sqrt{2x^2+3xy+4y^2}+\sqrt{2y^2+3xy+4x^2}=$\\
$\sqrt{\left[\sqrt{2}\left(x+\dfrac{3}{4}y\right)\right]^2+\left(\sqrt{\dfrac{23}{8}}y\right)^2}+\sqrt{\left[\sqrt{2}\left(y+\dfrac{3}{4}x\right)^2\right]+\left(\sqrt{\dfrac{23}{8}}x\right)^2}\ge 3|x+y|=3(x+y)$\\
Dấu bằng xảy ra khi $x=y\ge 0$. Đặt $x+y=t$ ta có $\va{t^2-t\ge 0 \\ t\ge 0}\iff \hay{t=0 \\ t\ge 1}~(*)$\\
Ta có: $P=2t^3+2t^2-xy(6t+5)+\sqrt{x^2+1}+\sqrt{y^2+1}\\
\ge 2t^3+2t^2-\dfrac{t^2}{4}(6t+5)+\sqrt{t^2+4}$\\
$\iff 4P\ge 2t^3+3t^2+4\sqrt{t^2+4}=f(t)$\\
Xét hàm số $f(t)=2t^3+3t^2+4\sqrt{t^2+4}$ trên $(*)$, $f'(t)\ge 0$ với mọi $t$ thỏa mãn $(*)$ suy ra \\
\cl{$f(t)\ge \left \{ f(0);f(1)\right \}=8$}\\
Vậy $\min P=2\iff x=y=0$.

\section{Lời giải Bộ 38 (THPT Thường Xuân 3, Thanh Hóa)}

\deda{
1. Trong mặt phẳng với hệ tọa độ $Oxy$, cho tam giác $ABC$ có đỉnh $A(3;-4)$. Phương trình đường trung trực cạnh $BC$, đường trung tuyến xuất phát từ $C$ lần lượt là $x+y-1=0$ và $3x-y-9=0$. Tìm tọa độ các đỉnh $B,C$ của tam giác $ABC$.\\
2. Trong mặt phẳng với hệ tọa độ $Oxy$, cho đường tròn $(C)$ có phương trình $x^2+y^2+2x-4y-8=0$ và đường thẳng $(\Delta)$ có phương trình: $2x-3y-1=0$. Chứng minh rằng $(\Delta)$ luôn cắt $(C)$ tại hai điểm phân biệt $A,B$. Tìm tọa độ điểm $M$ trên đường tròn $(C)$ sao cho diện tích tam giác $ABM$ lớn nhất.
}
\lgi 
\vspace{5cm}

\deda{
Tìm tất cả các giá trị thực của tham số $m$ để hệ sau có nghiệm thực:\\ \cl{$\begin{cases} x^2+\dfrac{4x^2}{(x+2)^2}\ge 5\\ x^4+8x^2+16mx+16m^2+32m+16=0 \end{cases}$.} 
}
\hda 
Với $x\neq -2$, PT $(1)$ tương đương\\
\cl{$\left(x-\dfrac{2x}{x+2}\right)^2+\dfrac{4x^2}{x+2}-5\ge 0\iff \left(\dfrac{x^2}{x+2}\right)^2+4.\dfrac{x^2}{x+2}-5=0\iff \hay{x\ge 2 \\ -2\neq x\le -1}$}\\
Giả sử $x_0$ là nghiệm của PT $(2)$.
Khi đó $16m^2+16(x_0+2)m+x_0^4+8x_0^2+16=0$ phải có nghiệm $m$\\
Do đó $\Delta '\ge 0\iff 0\le x_0\le 2$.
Nên hệ có nghiệm khi $(2)$ có nghiệm $x=2$\\
Thay $x=2\Rightarrow m=-2$.
Vậy $m=-2$.
\deda{
Tìm giá trị lớn nhất, giá trị nhỏ nhất của biểu thức\\
\cl{$P=\dfrac{\sqrt{5-4x}-\sqrt{1+a}}{\sqrt{5-4a}+2\sqrt{1+a}+6}$}
trong đó $a$ là tham số thực và $-1\le a\le \dfrac{5}{4}.$
}
\hda 
Đặt $A=\sqrt{5-4a} ;B=\sqrt{1+a}$ thì $A^2+4B^2=0; A,B\ge 0$\\
Do đó tồn tại $x\in \left[0;\dfrac{\pi}{2}\right]$: $A=3\sin x;2B=3\cos x$. Khi đó\\
\cl{$P=\dfrac{3\sin x-\dfrac{3}{2}\cos x}{3\sin x+3\cos x+6}=\dfrac{2\sin x-\cos x}{2\sin x+2\cos x+4}=f(x)$}\\
Ta có $f'(x)>0$ nên hàm đồng biến trên đoạn $x\in \left[0;\dfrac{\pi}{2}\right]$. Do đó \\
\cl{$Min f(x)=f(0)=-\dfrac{1}{6}; Max f(x)=f\left(\dfrac{\pi}{2}\right)=\dfrac{1}{3}$}\\
Vậy $minP=-\dfrac{1}{6}\iff a=\dfrac{5}{4}$, $maxP=\dfrac{1}{3}\iff a=-1$.

\section{Lời giải Bộ 39 (THPT Tĩnh Gia II)}

\deda{
Trong mặt phẳng với hệ tọa độ $Oxy$, cho đường tròn $(C): x^2+y^2=5$ tâm $O$, đường thẳng $(d):3x-y-2=0$. Tìm tọa độ các điểm $A,B$ trên $(d)$ sao cho $OA=\dfrac{\sqrt{10}}{5}$ và đoạn $OB$ cắt $(C)$ tại $K$ sao cho $KA=KB$.
}
\lgi 
\vspace{5cm}

\deda{
Giải hệ phương trình \\ \cl{$\begin{cases} \sqrt{x^2+2x+5}-\sqrt{y^2-2y+5}=y-3x-3\\y^2-3y+3=x^2-x  \end{cases}$.} 
}
\hda 
PT $(2)\iff y-3x-3=y^2-x^2-2y-2x$\\
Thay vào PT $(1)$ ta có\\
$\sqrt{(x+1)^2+}-\sqrt{(y-1)^2+4}=y^2-x^2-2y-2x\\
\iff \sqrt{(x+1)^2+4}+(x+1)^2=\sqrt{(y-1)^2+4}+(y-1)^2~(*)$\\
Xét hàm số $f(t)=\sqrt{t+4}+t$ trên $[0;+\infty)$, $f'(t)>0~\forall t\ge 0$\\
$\Rightarrow f(t)$ đồng biến trên $[0;+\infty)$\\
Vậy $(*)\iff (x+1)^2=(y-1)^2\iff \hay{x=y-2 \\ x=-y}$\\
Đáp án: $(x;y)=\left(-\dfrac{1}{2};\dfrac{3}{2}\right);\left(-\dfrac{3}{4};\dfrac{3}{4}\right)$.

\deda{
Cho các số thực dương $a,b,c$. Chứng minh rằng:\\
\cl{$\dfrac{\sqrt{a+b+c}+\sqrt{a}}{b+c}+\dfrac{\sqrt{a+b+c}+\sqrt{b}}{c+a}+\dfrac{\sqrt{a+b+c}+\sqrt{c}}{a+b}\ge \dfrac{9+3\sqrt{3}}{2\sqrt{a+b+c}}$.}
}
\hda 
ĐPCM$\iff \dfrac{1+\sqrt{\dfrac{a}{a+b+c}}}{\dfrac{b}{a+b+c}+\dfrac{c}{a+b+c}}+\dfrac{1+\sqrt{\dfrac{b}{a+b+c}}}{\dfrac{c}{a+b+c}+\dfrac{a}{a+b+c}}+\dfrac{1+\sqrt{\dfrac{c}{a+b+c}}}{\dfrac{a}{a+b+c}+\dfrac{b}{a+b+c}}\ge \dfrac{9+3\sqrt{3}}{2}$\\
Đặt $x=\dfrac{a}{a+b+c};y=\dfrac{b}{a+b+c};z=\dfrac{c}{a+b+c}$ ta có $x,y,z>0$ và $x+y+z=1$\\
Khi đó ĐPCM$\iff \dfrac{1+\sqrt{x}}{1-x}+\dfrac{1+\sqrt{y}}{1-y}+\dfrac{1+\sqrt{z}}{1-z}\ge \dfrac{9+3\sqrt{3}}{2}$\\
Ta có: $\dfrac{1}{1-x}+\dfrac{1}{1-y}+\dfrac{1}{1-z}\ge \dfrac{9}{3-x-y-z}=\dfrac{9}{2}$\\
Ta cm: $\sqrt{x}{1-x}+\dfrac{\sqrt{y}}{1-y}+\dfrac{\sqrt{z}}{1-z}\ge \dfrac{3\sqrt{3}}{2}$ dấu bằng xảy ra khi $x=y=z=\dfrac{1}{3}$\\
Xét hàm $f(x)=\sqrt{x}(1-x)$ với $0<x<1$ ta có $f'(x)=0\iff x=\dfrac{1}{3}$\\
$\Rightarrow 0<f(x)<\dfrac{2}{3\sqrt{3}}$ dấu bằng xảy ra khi $x=\dfrac{1}{3}$\\
Vậy ta có: $\dfrac{\sqrt{x}}{1-x}=\dfrac{x}{(1-x)\sqrt{x}}\ge \dfrac{3\sqrt{3}x}{2}$\\
Tương tự ta có $\sqrt{x}{1-x}+\dfrac{\sqrt{y}}{1-y}+\dfrac{\sqrt{z}}{1-z}\ge \dfrac{3\sqrt{3}}{2}$
Ta có đpcm, dấu bằng khi $x=y=z=\dfrac{1}{3}$.

\section{Lời giải Bộ 40 (THPT Triệu Sơn 3, Thanh Hóa)}

\deda{Trong mặt phẳng với hệ tọa độ $Oxy$ cho hình chữ nhật $ABCD$ có điểm $H(1;2)$ là hình chiếu vuông góc của $A$ lên $BD$. Điểm $M\left(\dfrac{9}{2};3\right)$ là trung điểm của cạnh $BC$, phương trình đường trung tuyến kẻ từ $A$ của $\Delta ADH$ là $d:4x+y-4=0$. Viết phương trình cạnh $BC$.
}
\lgi 
\vspace{5cm}

\deda{
Giải hệ phương trình \\ \cl{$\begin{cases} x\sqrt{x^2+y}+y=\sqrt{x^4+x^3}+x\\ x+\sqrt{y}+\sqrt{x-1}+\sqrt{y(x-1)}=\dfrac{9}{2}  \end{cases}$.} 
}
\hda 
ĐK: $x\ge 1,y\ge 0$\\
Thực hiện nhân liên hợp ta có PT $(1)\iff (x-y)\left(\sqrt{x^2+y}+\sqrt{x^2+x}-x\right)=0\\
\iff x=y$\\
Thay vào PT $(2)$ ta có $x+\sqrt{x}+\sqrt{x-1}+\sqrt{x(x-1)}=\dfrac{9}{2}$\\
Đặt $t=\sqrt{x}+\sqrt{x-1}, t\ge 0$ thì pt trở thành $t^2+1+2t=9$
$\iff t=2\\ \iff \sqrt{x}+\sqrt{x-1}=2\iff 2\sqrt{x(x-1)}=5-2x\\ \iff \va{x\le \dfrac{5}{2} \\ 4x^2-4x=25-20x+4x^2}\iff x=\dfrac{25}{16}$\\
Vậy hệ có nghiệm duy nhất $(x;y)=\left (\dfrac{25}{16};\dfrac{25}{16}\right)$.

\deda{
Cho $a,b,c$ thuộc khoảng $(0;1)$ thỏa mãn $\left (\dfrac{1}{a}-1\right)\left (\dfrac{1}{b}-1\right)\left (\dfrac{1}{c}-1\right)=1$. Tìm GTNN của biểu thức $P=a^2+b^2+c^2$.
}
\hda  
Từ $GT\iff ab+bc+ca=a+b+c-1+2abc$\\
Ta có: \\
$P=(a+b+c)^2-2(ab+bc+ca)=(a+b+c)^2-2(a+b+c-1)-4abc\\
\ge t^2-2t+2-\dfrac{4}{27}t^3$ với $t=a+b+c, 0<t<3$\\
Khảo sát hàm số trên tìm ra $minP=\dfrac{3}{4}$ khi $t=\dfrac{3}{2}$ hay $a=b=c=\dfrac{1}{2}$.

\section{Lời giải Bộ 41 (Trung tâm dạy thêm văn hóa Lê Hồng Phong)}

\deda{Trong mặt phẳng $Oxy$ cho hình thang $ABCD$ có đáy lớn $CD=3AB$, $C(-3;-3)$, trung điểm của $AD$ là $M(3;1)$. Tìm tọa độ đỉnh $B$ biết $S_{BCD}=18, AB=\sqrt{10}$ và đỉnh $D$ có hoành độ nguyên dương.}
\lgi 
\vspace{5cm}

\deda{
Giải hệ phương trình \\ \cl{$\begin{cases} x-y\sqrt{2-x}+2y^2=2\\ 2\left(\sqrt{x+2}-4y\right)+8\sqrt{y}\sqrt{xy+2y}=34-15x \end{cases}$.} 
}
\hda 
ĐK: $-2\le x\le 2;y\ge 0$\\
PT $(1)\iff (2-x)+\sqrt{2-x}.y-2y^2=0\iff \hay{\sqrt{2-x}=y \\ \sqrt{2-x}=-2y}$\\
+ Với $\sqrt{2-x}=y$, thay vào PT $(2)$ ta có $2\left ( \sqrt{x+2}-4\sqrt{2-x}\right)+8\sqrt{4-x^2}=34-15x$\\
Đặt $t=\sqrt{x+2}-4\sqrt{2-x}$, pt trở thành $2t=t^2\iff \hay{t=0 \\ t=2}$\\
$\hay{\sqrt{x+2}=4\sqrt{2-x} \\ \sqrt{x+2}=4\sqrt{2-x}+2}$\\
$\iff \hay{x=\dfrac{30}{17}\Rightarrow y=\dfrac{2\sqrt{17}}{17} \\ x=2\Rightarrow y=0}$\\
+ Với $\sqrt{2-x}=-2y\ge 0\Rightarrow y\le 0$, mà $y\ge 0\Rightarrow y=0\Rightarrow x=2$\\
Vậy hệ đã cho có 2 nghiệm $(x;y)=(2;0),\left( \dfrac{30}{17}; \dfrac{2\sqrt{17}}{17}\right)$.

\deda{
Cho $x,y$ là các số không âm thỏa $x^2+y^2=2$. Tìm giá trị lớn nhất và nhỏ nhất của:\\
\cl{$P=5(x^5+y^5)+x^2y^2\left(5\sqrt{2xy+2}-4xy+12\right)$.}
}
\hda 
Với $0\le x,y\le \sqrt{2}\Rightarrow \va{x^2(x-\sqrt{2})\le 0 \\ y^2(y-\sqrt{2})\le 0}\Rightarrow x^3+y^3\le 2\sqrt{2}$\\
Có: $(x+y)^2\le 2(x^2+y^2)=4\Rightarrow x+y\le 2$\\
$\Rightarrow 2(x^3+y^3)\ge (x+y)(x^3+y^3)\ge (x^2+y^2)^2=4$\\
$\Rightarrow x^3+y^3\ge 2$\\
Đặt $t=x^3+y^3, t\in [2;2\sqrt{2}]$\\
Ta có:\\
+ $2^3=(x^2+y^2)^3=(x^3+y^3)^2-2x^3y^3+6x^2y^2$\\
$\Rightarrow 2x^3y^3-6x^2y^2=t^2-8$\\
+ $2(x^3+y^3)=(x^3+y^3)(x^2+y^2)=x^5+y^5+x^2y^2(x+y)$\\
$\Rightarrow x^5+y^5+x^2y^2(x+y)=2t$\\
Vậy $P=-4x^3y^3+12x^2y^2+5(x^5+y^5)+5x^2y^2\sqrt{2+2xy}\\
=-2t^2+10t+16=f(t)$\\
Có: $f'(t)=-4t+10=0\iff t=\dfrac{5}{2}$ (thỏa mãn)\\
Vậy $minP=f(2)=28$ và $maxP=f\left(\dfrac{5}{2}\right)=\dfrac{57}{2}$.

\section{Lời giải Bộ 42 (THPT chuyên Vĩnh Phúc, Vĩnh Phúc lần 2)}

\deda{
Trong mặt phẳng với hệ tọa độ $Oxy$, cho đường tròn $(C): x^2+y^2-2x-4y-4=0$ tâm $I$ và điểm $M(3;2)$. Viết phương trình đường thẳng $\Delta$ đi qua $M$, $\Delta$ cắt $(C)$ tại hai điểm phân biệt $A,B$ sao cho diện tích tam giác $IAB$ lớn nhất.
}
\lgi 
\vspace{5cm}

\deda{
Giải hệ phương trình \\ \cl{$\begin{cases} x^4-2x=y^4-y\\ (x^2-y^2)^3=3   \end{cases}$.} 
}
\hda 
Đặt $x+y=a;x-y=b,3=c^3$, từ PT $(2)$ ta có $ab=c$\\
Từ $x=\dfrac{a+b}{2};y=\dfrac{a-b}{2}\Rightarrow x^4-y^4=\dfrac{ab}{2}(a^2+b^2)$ và\\
$(2x-y)=a+b-\dfrac{a-b}{2}=\dfrac{a+3b}{2}=\dfrac{a+bc^3}{2}$\\
Ta có hệ $\va{c(a^2+b^2)=a+bc^3 \\ ab=c}$\\
$\Rightarrow ca^4+c^3=a^3+ac^4\\ 
\iff (ca-1)(a^3-c^3)=0\\
\iff \hay{a=\dfrac{1}{c} \\ a=c}$\\
+ Nếu $a=c\Rightarrow b=1$ thì $x=\dfrac{c+1}{2}=\dfrac{\sqrt[3]{3}+1}{2};y=\dfrac{\sqrt[3]{3}-1}{2}$\\
+ Nếu $a=\dfrac{1}{c}\Rightarrow b=c^2$ thì $x=\dfrac{2}{\sqrt[3]{3}},y=-\dfrac{1}{\sqrt[3]{3}}$\\
Vậy hệ đã cho có hai nghiệm là $(x;y)=\left(\dfrac{\sqrt[3]{3}+1}{2};\dfrac{\sqrt[3]{3}-1}{2}\right);\left(\dfrac{2}{\sqrt[3]{3}}; -\dfrac{1}{\sqrt[3]{3}}\right)$.

\deda{
Cho các số $a,b,c$ không âm sao cho tổng hai số bất kì đều dương. Chứng minh rằng:\\
\cl{$\sqrt{\dfrac{a}{b+c}}+\sqrt{\dfrac{b}{c+a}}+\sqrt{\dfrac{c}{a+b}}+\dfrac{9\sqrt{ab+bc+ca}}{a+b+c}\ge 6$.}
}
\hda 
Giả sử $a\ge b\ge c$, khi đó $\sqrt{\dfrac{ab}{a+c}}+\sqrt{\dfrac{ac}{a+b}}\ge \sqrt{\dfrac{b.b}{b+c}}+\sqrt{\dfrac{c.c}{c+b}}=\sqrt{b+c}$
$\Rightarrow \sqrt{\dfrac{b}{a+c}}+\sqrt{\dfrac{c}{a+b}}\ge \sqrt{\dfrac{b+c}{a}}$\\
Đặt $t=b+c$ thì $VT\ge \sqrt{\dfrac{a}{t}}+\sqrt{\dfrac{t}{a}}+\dfrac{9\sqrt{at}}{a+t}$.
Ta có $\sqrt{\dfrac{a}{t}}+\sqrt{\dfrac{t}{a}}+\dfrac{9\sqrt{at}}{a+t}=\dfrac{a+t}{\sqrt{at}}+\dfrac{9\sqrt{at}}{a+t}\ge 6$ (AM-GM).
Do đó $VT\ge 6$ (dpcm)\\
\textbf{Chú ý:} Đẳng thức xảy ra khi $a+t=3\sqrt{at}$ khi chẳng hạn $(a;b;c)=\left(\dfrac{7+3\sqrt{5}}{2};1;0\right)$ (HS có thể bỏ qua bước này).

\section{Lời giải Bộ 43 (THPT Đồng Lộc, Hà Tĩnh)}

\deda{ Trong mặt phẳng với hệ tọa độ $Oxy$, cho hình thang cân $ABCD$ có diện tích bằng $\dfrac{45}{2}$, đáy lớn $CD$ nằm trên đường thẳng $x-3y-3=0$. Biết hai đường chéo $AC,BD$ vuông góc với nhau tại $I(2;3)$. Viết phương trình đường thẳng chứa cạnh $BC$, biết điểm $C$ có hoành độ dương.}
\lgi 
\vspace{5cm}

\deda{
Giải hệ phương trình \\ \cl{$\begin{cases} y^3+6y^2+16y-3x+11=0\\ x^3+3x^2+x+3y+3=0   \end{cases}$.} 
}
\hda 
HPT$\iff \va{(y+2)^3+4y-3x+3=0 \\ (x+1)^3+3y-2x+2=0}$\\
Đặt $\va{a=y+2\\b=x+1}\iff \va{y=a-2\\x=b-1}$, thay vào hệ ta có\\
$\va{a^3+4a-3b=2\\b^3+3a-2b=2}$\\
Trừ theo vế ta được:\\
$(a-b)(a^2+ab+b^2+1)=0\iff a=b$\\
Thay $a=b$ vào PT $(1)$ ta có\\
$a^3+a-2=0\iff \hay{a=1\Rightarrow b=1\\a^2+a+2=0(VN)}$\\
Khi đó $x=0;y=-1$\\
Vậy $x=0;y=-1$.


\deda{
Cho $0<a,b,c<\dfrac{1}{2}$ thỏa mãn $a+2b+3c=2$. Chứng minh rằng:\\
\cl{$\dfrac{1}{a(4b+6c-3)}+\dfrac{2}{b(3c+a-1)}+\dfrac{9}{c(2a+4b-1)}\ge 54$.}
}

\hda 
\textbf{Bổ đề} Với $x>0$. Chứng minh $x^2(1-2x)\le \dfrac{1}{27}$\\
Ta có: $x^2(1-2x)=x.x.(1-2x)\le \left(\dfrac{x+x+1-2x}{3}\right)^3=\dfrac{1}{27}$(AM-GM)\\
+ Ta có:\\
$A=\dfrac{1}{a(1-2a)}+\dfrac{2}{b(1-2b)}+\dfrac{9}{c(3-6c)}\\
=\dfrac{a}{a^2(1-2a)}+\dfrac{2b}{b^2(1-2b)}+\dfrac{3c}{c^2(1-2c)}$\\
Áp dụng bổ đề trên ta có\\
$A\ge 27a+27.2b+27.3c=54$\\
Dấu bnawgf xảy ra khi $a=b=c=\dfrac{1}{3}$\\

\section{Lời giải Bộ 44 (THPT Hậu Lộc 2, Thanh Hóa)}

\deda{
Trong không gian tọa độ $Oxyz$ cho mặt phẳng $(P): x-y+z+2=0$ và điểm $A(1;-1;2)$. Tìm tọa độ điểm $A'$ đối xứng với điểm $A$ qua mặt phẳng $(P)$. Viết phương trình mặt cầu đường kính $AA'$.
}
\lgi 
\vspace{5cm}

\deda{
Giải hệ phương trình \\ \cl{$\begin{cases} (x+1)^2+y^2=2\left(1+\dfrac{1-y^2}{x}\right)\\ 4y^2=(y^2-x^3+3x-2)\left(\sqrt{2-x^2}+1\right)   \end{cases}$.}
}
\hda 
ĐK: $x\ne 0,-\sqrt{2}\le x\le \sqrt{2}$\\
PT $(1)\iff x(x+1)^2+xy^2=2(x+1-y^2)\\
\iff (x+2)(x^2+y^2-1)=0\\
\iff \hay{x=-2 (L) \\ x^2+y^2=1}$\\
+ Với $x^2+y^2=1\iff x^2=1-y^2$, thay vào PT $(2)$ ta có:\\
$4y^2=(y^2-x^3+3x-2)\left(\sqrt{y^2+1}+1\right)\\
\iff 4\left(\sqrt{y^2+1}+1\right)\left(\sqrt{y^2+1}-1\right)=(y^2-x^3+3x-2)\left(\sqrt{y^2+1}+1\right)\\
\iff x^3-3x-2=y^2-4\sqrt{y^2+1}$\\
+ Do $x^2+y^2=1\Rightarrow \va{-1\le x\le 1 \\ -1\le y\le 1}$\\
+ Xét hàm $f(x)=x^3-3x-2$ trên đoạn $[-1;1]$\\
Có $f'(x)=3x^2-3=0\iff x=\pm 1$\\
Do hàm $f(x)$ liên tục trên đoạn $[-1;1]$ và $f(-1)=0;f(1)=-4$ nên\\
\cl{$minf(x)=-4$~~~~, $maxf(x)=0$}\\
+ Xét hàm số $g(y)=y^2-4\sqrt{y^2+1}$ trên đoạn $[-1;1]$ ta được \\
\cl{$min g(y)=1-4\sqrt{2}~~~~~,maxg(y)=-4$}\\
Vậy ta có $f(x)=g(y)=-4\iff \va{x=1 \\ y=0}$ (thỏa mãn)\\
Vậy hệ có nghiệm duy nhất $(x;y)=(1;0)$.
\deda{
Cho các số thực dương $x,y,z$ thỏa mãn $xy\ge 1, z\ge 1$. Tìm giá trị nhỏ nhất của biểu thức:\\
\cl{$P=\dfrac{x}{y+1}+\dfrac{y}{x+1}+\dfrac{z^3+2}{3(xy+1)}$}
}
\hda 
\textbf{Bổ đề:} Với $x,y>0$ thỏa mãn $xy\ge 1$ ta có $\dfrac{1}{x+1}+\dfrac{1}{y+1}\ge \dfrac{2}{1+\sqrt{xy}}~(1)$\\
Biến đổi tương đương ta được $(1)\iff (\sqrt{xy}-1)(\sqrt{x}-\sqrt{y})^2\ge 0$ luôn đúng do $xy\ge 1$ (dpcm)\\
+ Theo BĐT AM-GM ta có $z^3+2=z^3+1+1\ge 3z\ge 3$\\
$\Rightarrow P\ge \dfrac{x}{y+1}+1+\dfrac{y}{x+1}+1+\dfrac{1}{xy+1}-2\\
=(x+y+1)\left(\dfrac{1}{x+1}+\dfrac{1}{y+1}\right)+\dfrac{1}{xy+1}-2\ge (2\sqrt{xy}+1).\dfrac{2}{1+\sqrt{xy}}+\dfrac{1}{xy+1}-2$\\
+ Đặt $t=\sqrt{xy}, t\ge 1$ ta có:
$P\ge P(t)=(2t+1).\dfrac{2}{t+1}+\dfrac{1}{t^2+1}-2=\dfrac{2t}{t+1}+\dfrac{1}{t^2+1}$\\
Ta có $P'(t)\ge 0 \forall t\ge 1$
$\Rightarrow P(t)\ge P(1)=\dfrac{3}{2}\forall t\ge 1$\\
Vậy $minP=\dfrac{3}{2}\iff x=y=z=1$.

\section{Lời giải Bộ 45 (THPT Nguyễn Trung Thiên)}

\deda{
Trong mặt phẳng với hệ tọa độ $Oxy$, cho tam giác $ABC$ có trung điểm cạnh $BC$ là $M(3;-1)$. Điểm $E(-1;-3)$ nằm trên đường thẳng $\Delta$ chứa đường cao qua đỉnh $B$. Đường thẳng $AC$ qua $F(1;3)$. Tìm tọa độ các đỉnh của $\Delta ABC$ biết đường tròn ngoại tiếp $\Delta ABC$ có đường kính $AD$ với $D(4;-2)$.
}
\lgi 
\vspace{5cm}

\deda{
Giải phương trình: $(x+2)\left(\sqrt{x^2+4x+7}+1\right)+x\left(\sqrt{x^2+3}+1\right)=0$.
}
\hda 
PT$\iff(x+2)\left(\sqrt{(x+2)^2+3}+1\right) =-x\left(\sqrt{x^2+3}+1\right)$\\
+ Đặt $u=x+2;v=-x$.
Xét hàm $f(t)=t\left(\sqrt{t^2+3}+1\right)$, pt trở thành $f(u)=f(v)$\\
Vì $f'(t)>0\forall t$ nên $f(u)=f(v)\iff u=v$.
PT$\iff x+2=-x\iff x=-1$\\
Vậy pt có nghiệm duy nhất $x=-1$.

\deda{
Cho $x,y,z$ là ba số thực dương thỏa mãn $x^2+y^2+z^2=3.$ Tìm giá trị nhỏ nhất của biểu thức sau:\\
\cl{$P=\dfrac{x}{(y+z)^2}+\dfrac{y}{(x+z)^2}+\dfrac{z}{(x+y)^2}$.}
}
\hda 
Giả thiết suy ra $0<x,y,z<\sqrt{3}$, ta có $(y+z)^2\le 2(y^2+z^2)=2(3-x^2)$\\
$\Rightarrow \dfrac{x}{(y+z)^2}\ge \dfrac{x}{2(3-x^2)}$, mặt khác $\dfrac{x}{2(3-x^2)}\ge \dfrac{1}{4}x^2\iff (x-1)^2(x+2)\ge 0$ (luôn đúng)\\
Tương tự ta có $P\ge \dfrac{1}{4}(x^2+y^2+z^2)=\dfrac{3}{4}$\\
Vậy $minP=\dfrac{3}{4}\iff x=y=z=1$.
